\documentclass[10pt,a4paper]{amsart}
\usepackage[margin=19mm]{geometry}
\usepackage{amsmath,amssymb,mathtools}
\usepackage[scr=rsfs,cal=euler]{mathalfa}
\usepackage{xfrac}
\usepackage[unicode,hidelinks,bookmarks=false]{hyperref}
\newcommand{\ie}{i.e.\ }
\newcommand{\cf}{cf.\ }
\newcommand{\resp}{respectively\ }
\newcommand{\Subsection}{\subsection}
\providecommand{\nobreakdash}{\nobreak\hbox{-}}
\setlength{\emergencystretch}{2em}
\multlinegap0pt
\usepackage{amsthm}
\usepackage{mathtools}
\usepackage{amsmath,amssymb}
\usepackage{bm}
\usepackage{placeins}
\usepackage{graphicx}
\newtheorem{assumption}{Assumption}
\newtheorem{theorem}{Theorem}
\newtheorem{definition}{Definition}
\newtheorem{lemma}{Lemma}
\newtheorem{corollary}{Corollary}
\title[PowerStep: memory-efficient adaptive optimization]{PowerStep: Memory-Efficient Adaptive Optimization via $\ell_p$-Norm Steepest Descent: the convergence theorem with its explicit constants (Theorem 1)}
\author[Y. Lu, D. Fan, S. Zhang, Y. Tian]{Yao Lu, Dengdong Fan, Shixun Zhang, Yonghong Tian}
\date{Source: arXiv:2605.10335v2 (CC BY 4.0)}
\begin{document}
\maketitle
\vspace{1em}
\noindent\textit{Source and licence.} PowerStep: Memory-Efficient Adaptive Optimization via $\ell_p$-Norm Steepest Descent. Version arXiv:2605.10335v2 (CC BY 4.0), distributed under the Creative Commons Attribution 4.0 International licence, http://creativecommons.org/licenses/by/4.0/, as recorded in the arXiv OAI licence metadata for this exact version and shown on the abstract page https://arxiv.org/abs/2605.10335v2. This item is an editorial re-presentation of the paper's convergence theorem together with the paper's own proof of it and the proof-local material it uses. A full rights notice is delivered at licenses/arxiv-2605.10335-CC-BY-4.0.txt.
\noindent\textit{Editorial scope.} Counted unit: the paper's convergence theorem for its $\ell_p$-norm steepest-descent method, printed as Theorem 1 (source0.tex lines 163--172), delivered together with the paper's complete original proof of it (source0.tex lines 852--918, the appendix subsection that the paper devotes to that proof) as the single primary proof body. No mathematics is written, repaired or re-derived: statement and proof are the paper's own text line for line, in the paper's own notation ($f$, $\bm{\theta}_t$, $\bar{\mathbf{g}}_t$, $\mathbf{u}_t$, $L_p$, $G$, $\sigma$, $\gamma$, $\beta$, $q$, $h$, $R_T$, $C_\beta$). Required context retained uncounted, because the counted proof is the paper's own assembly over it: the two standing assumptions (source0.tex lines 152--160); the appendix opening with the buffer decomposition display (lines 715--724); the explicit-constants and dimension-conversion displays (lines 726--753); and the four transform and moment lemmas with their proofs (lines 757--849). The paper states its proof of the counted theorem as an appendix subsection rather than as a proof environment; this package delivers exactly that subsection as the primary proof body, and the delivered evidence says so. Every cross-reference inside the retained spans resolves inside the delivered document, and no citation command occurs in any retained span, so the delivered document carries no bibliography and no reference or citation command of the delivered text was rewritten or adapted. Printed numbers are the paper's own: the wrapper restores the paper's per-kind counters and its global equation counter before the retained material, so the retained assumptions print as Assumptions 1 and 2, the counted statement as Theorem 1, the retained lemmas as Lemmas 1 to 4, and the numbered displays of the retained spans carry the paper's own equation numbers (10), (11) and (21) to (43). Omitted from this package and not redistributed: the paper's preamble and front matter as such (of which only the title and the author names are reproduced in the wrapper); the introduction, the method section and the experiments; the corollaries and the remaining appendix material; and the paper's whole bibliography with its bib data file. Printed slips kept literal, among them the paper's own wording and its own choice of symbols. Layout and class: the paper's own class is the standard LaTeX article class with the publisher's conference style file for the submission, which is neither shipped with this package nor used to build it; the delivered wrapper is the lane's pinned amsart editorial wrapper at 10pt on A4, to which this package adds the paper's own theorem-like declarations with their separate counters and the shorthand macros its retained text uses, all copied verbatim from the paper's source and its own local style file. No publisher class file, no publisher style file and no upstream script is redistributed. Assets: no retained span of this package contains a figure environment or a graphics-inclusion command, no image or artwork file exists in or is redistributed with this package, and the manifest and the delivery evidence both carry an empty asset-dependency list. A full rights notice is delivered at licenses/arxiv-2605.10335-CC-BY-4.0.txt.
\setcounter{equation}{9}
\section{Convergence Analysis}
\label{sec:convergence}

\begin{assumption}[Smoothness and boundedness]\label{assum:regularity}
For constants $L_p,G>0$, the continuously differentiable objective is bounded below by $f^*$, satisfies $\|\nabla f(\mathbf{x})\|_q\le G$, and has a Lipschitz gradient in the dual pair:
\begin{equation}
\|\nabla f(\mathbf{x})-\nabla f(\mathbf{y})\|_q\le L_p\|\mathbf{x}-\mathbf{y}\|_p.
\end{equation}
\end{assumption}
\begin{assumption}[Stochastic oracle]\label{assum:oracle}
With $\mathcal F_{t-1}$ the history before sampling $\mathbf{g}_t$, write $\bar{\mathbf{g}}_t=\nabla f(\bm{\theta}_{t-1})$ and $\bm{\xi}_t=\mathbf{g}_t-\bar{\mathbf{g}}_t$. We assume $\mathbb E[\|\bm{\xi}_t\|_q^q\mid\mathcal F_{t-1}]\le\sigma^q$; the oracle may be biased.
\end{assumption}

\begin{theorem}[Finite-horizon stationarity bound]\label{thm:convergence}
Under Assumptions~\ref{assum:regularity}--\ref{assum:oracle}, fix a horizon $T$, use $\eta_t=\eta=c/\sqrt T$ with $c>0$, and initialize $\mathbf{m}_0=\mathbf{0}$. Let $\nu_t=\mathbb E\|\bm{\xi}_t\|_q^q$. There are explicit constants $C_1,C_2\ge0$ and $C_\beta>0$, independent of $T$, such that
\begin{equation}\label{eq:corrected-rate}
\frac1T\sum_{t=1}^T\mathbb E\|\bar{\mathbf{g}}_t\|_q^q
\le \underbrace{\frac{C_1}{\sqrt T}+\frac{C_2}{T}}_{R_T}
+\frac{C_\beta}{T}\sum_{t=1}^T\nu_t.
\end{equation}
Here $C_1,C_2$ depend on $L_p,G,\sigma,\gamma,\beta,c$ and the initial objective gap; $C_\beta$ depends only on $\beta$. Appendix~\ref{app:proofs} gives their formulas.
Consequently, $\min_{t\le T}\mathbb E\|\bar{\mathbf{g}}_t\|_2^2\le G^{1-\beta}(R_T+C_\beta\sigma^q)$.
\end{theorem}

\setcounter{equation}{20}
\FloatBarrier
\section{Complete Convergence Proofs}
\label{app:proofs}
Throughout this appendix, $0<\beta\le1$, $q=1+\beta$, $p=q/\beta$, $h=1-\gamma$, and the assumptions are those of Section~\ref{sec:convergence}. All vector norms below are explicitly indicated. The update is exact, with no weight decay, clipping, or quantization. Expectations refer to the stochastic oracle, and the horizon-dependent step size is constant during each run. Write $\bar{\mathbf{g}}_t=\nabla f(\bm{\theta}_{t-1})$, $\mathbf{u}_t=\Phi_\beta(\mathbf{m}_t)$, and
\begin{equation}\label{eq:buffer-decomp}
\bar{\mathbf{G}}_t=\sum_{k=0}^{t-1}\gamma^k\bar{\mathbf{g}}_{t-k},\qquad
\mathbf{N}_t=\sum_{k=0}^{t-1}\gamma^k\bm{\xi}_{t-k},\qquad
\mathbf{m}_t=\bar{\mathbf{G}}_t+\mathbf{N}_t,\qquad \mathbf{r}_t=\bar{\mathbf{g}}_t-h\bar{\mathbf{G}}_t.
\end{equation}
Here $\bar{\mathbf{G}}_t$ is an accumulated true-gradient component, not a conditional expectation of $\mathbf{m}_t$ at the current iterate. Past noise is correlated with later iterates; no independence between $\mathbf{N}_t$ and $\bar{\mathbf{G}}_t$ is assumed.

\paragraph{Explicit constants in Theorem~\ref{thm:convergence}.}
Let $\Delta_0=f(\bm{\theta}_0)-f^*$. Define
\begin{align}
M&=\left(\frac{G+\sigma}{h}\right)^{2\beta},
&D&=\frac{L_p\sqrt M\,\gamma}{h},\nonumber\\
A&=2^{-q}h^{-\beta},
&B&=\sqrt M+\tfrac12h^{-\beta}(2G)^\beta,\label{eq:theory-constants}\\
K_\beta&=\frac{\beta}{2}\left(\frac{2^{2-\beta}}q\right)^{q/\beta}.\nonumber
\end{align}
Theorem~\ref{thm:convergence} uses
\begin{equation}
C_1=\frac1A\left[\frac{\Delta_0}{c}+c\left(\frac{L_pM}{2}+BD\right)\right],\qquad
C_2=\frac{BG\gamma}{Ah},\qquad C_\beta=2^qK_\beta.
\end{equation}
Thus $R_T=C_1/\sqrt T+C_2/T$ retains the initialization term and all parameter dependencies.

\paragraph{Proof idea and scope.}
Write $\mathbf{m}_t=\bar{\mathbf{G}}_t+\mathbf{N}_t$, where $\bar{\mathbf{G}}_t=\sum_{k=0}^{t-1}\gamma^k\bar{\mathbf{g}}_{t-k}$ and $\mathbf{N}_t=\sum_{k=0}^{t-1}\gamma^k\bm{\xi}_{t-k}$. Align the update through $\bar{\mathbf{G}}_t$ and control $\bar{\mathbf{g}}_t-h\bar{\mathbf{G}}_t$ in mean square. This makes the drift contribution linear in $\eta$, rather than $\eta^\beta$. A coordinatewise H\"older bound and Young's inequality control the noise without removing its residual. The following lemmas account for initialization and give explicit constants. The constant horizon-dependent step avoids the logarithm in $\sum_t(c/\sqrt t)^2$; it differs from the experimental warmup/cosine schedule.

When $\sigma=0$, the bound gives $O(T^{-1/2})$ stationarity. A vanishing average noise moment also removes the residual. The corollaries below state sufficient conditions for growing batches and a small relative-noise condition. These are additional assumptions, not properties established for the language-model experiments. We make no minimax-optimality claim: a neighborhood bound is weaker than exact stationarity, and increasing the batch changes the number of oracle calls.

\paragraph{Dimension and the choice of $\beta$.}
There is no explicit $d$ in~\eqref{eq:corrected-rate} when $L_p,G,\sigma$ are fixed in the native norms. This does not imply dimension independence at fixed Euclidean constants. If those constants are $L_2,G_2,\sigma_2$, valid choices are
\begin{equation}\label{eq:dimension-conversion}
L_p=d^{(1-\beta)/(1+\beta)}L_2,\quad
G=d^{1/q-1/2}G_2,\quad
\sigma^q=d^{1-q/2}\sigma_2^q.
\end{equation}

\begin{lemma}[Dual geometry and scalar perturbations]\label{lemma:holder}
For $\mathbf{x},\mathbf{y},\mathbf{m}\in\mathbb R^d$,
\begin{equation}
\langle \mathbf{m},\Phi_\beta(\mathbf{m})\rangle=\|\mathbf{m}\|_q^q,\quad
\|\Phi_\beta(\mathbf{m})\|_p=\|\mathbf{m}\|_q^\beta,\quad
\|\Phi_\beta(\mathbf{x})-\Phi_\beta(\mathbf{y})\|_p\le2^{1-\beta}\|\mathbf{x}-\mathbf{y}\|_q^\beta.
\end{equation}
Moreover, for real $a,n$,
\begin{equation}\label{eq:young-explicit}
2^{1-\beta}|a||n|^\beta\le\tfrac12|a|^q+K_\beta|n|^q,
\qquad K_\beta=\frac\beta2\left(\frac{2^{2-\beta}}q\right)^{q/\beta}.
\end{equation}
\end{lemma}
\begin{proof}
The identities follow coordinatewise from $\beta p=q$. For same-sign scalars, subadditivity of $z^\beta$ gives $|\Phi_\beta(x)-\Phi_\beta(y)|\le|x-y|^\beta$. For opposite signs, concavity gives $|x|^\beta+|y|^\beta\le2^{1-\beta}(|x|+|y|)^\beta$. Raising the scalar bound to $p$, summing, and taking the $p$-th root proves the vector bound. Finally maximize $2^{1-\beta}a n^\beta-a^q/2$ over $a\ge0$ for fixed $n\ge0$. For $n>0$ the maximizer is $a=(2^{2-\beta}/q)^{1/\beta}n$, giving~\eqref{eq:young-explicit}; the zero cases are immediate.
\end{proof}

\begin{lemma}[Uniform update moment and accumulated noise]\label{lemma:momentum_bound}
Let $\nu_t=\mathbb E\|\bm{\xi}_t\|_q^q\le\sigma^q$. Then
\begin{align}
\mathbb E\|\mathbf{u}_t\|_p^2&\le M=\left(\frac{G+\sigma}{h}\right)^{2\beta},\label{eq:update-moment}\\
\mathbb E\|\mathbf{N}_t\|_q^q&\le h^{1-q}\sum_{k=0}^{t-1}\gamma^k\nu_{t-k},\qquad
\sum_{t=1}^T\mathbb E\|\mathbf{N}_t\|_q^q\le h^{-q}\sum_{t=1}^T\nu_t.\label{eq:noise-moment}
\end{align}
\end{lemma}
\begin{proof}
Minkowski's inequality in $L^q$ gives $(\mathbb E\|\mathbf{g}_t\|_q^q)^{1/q}\le G+\sigma$ and
$(\mathbb E\|\mathbf{m}_t\|_q^q)^{1/q}\le\sum_{k=0}^{t-1}\gamma^k(G+\sigma)\le(G+\sigma)/h$.
Since $2\beta/q\le1$, Jensen yields
$\mathbb E\|\mathbf{u}_t\|_p^2=\mathbb E(\|\mathbf{m}_t\|_q^q)^{2\beta/q}\le(\mathbb E\|\mathbf{m}_t\|_q^q)^{2\beta/q}\le M$.
This does not require a second moment of the untransformed noise.
For nonnegative weights $w_k$, convexity gives
$\|\sum_k w_k\mathbf{x}_k\|_q^q\le(\sum_k w_k)^{q-1}\sum_k w_k\|\mathbf{x}_k\|_q^q$.
Set $w_k=\gamma^k$ and take expectations. Summing and interchanging the finite sums bounds the remaining geometric sum by $1/h$. No orthogonality or independence of the noise terms is used. Neither this estimate nor the subsequent alignment argument requires a zero conditional mean; oracle bias is included in the noise moment and its residual.
\end{proof}

\subsection{Drift and Alignment}
\begin{lemma}[Initialization and gradient drift]\label{lemma:drift}
With constant step size $\eta$, let $D=L_p\sqrt M\gamma/h$. Then
\begin{equation}\label{eq:drift-rms}
\left(\mathbb E\|\mathbf{r}_t\|_q^2\right)^{1/2}\le G\gamma^t+D\eta=:R_t^{\rm drift},
\qquad \|\mathbf{r}_t\|_q\le2G\ \text{almost surely}.
\end{equation}
Throughout, $R_t^{\rm drift}$ denotes this deterministic RMS upper bound; it also bounds $\mathbb E\|\mathbf r_t\|_q$ by Cauchy--Schwarz. Consequently,
\begin{equation}\label{eq:drift-products}
\mathbb E[\|\mathbf{r}_t\|_q\|\mathbf{u}_t\|_p]\le\sqrt M R_t^{\rm drift},
\qquad \mathbb E\|\mathbf{r}_t\|_q^q\le(2G)^{q-1}R_t^{\rm drift}.
\end{equation}
\end{lemma}
\begin{proof}
From~\eqref{eq:buffer-decomp},
\begin{equation}
\mathbf{r}_t=\gamma^t\bar{\mathbf{g}}_t+h\sum_{k=1}^{t-1}\gamma^k(\bar{\mathbf{g}}_t-\bar{\mathbf{g}}_{t-k}).
\end{equation}
The first term is bounded by $G\gamma^t$. By the gradient Lipschitz assumption,
$\|\bar{\mathbf{g}}_t-\bar{\mathbf{g}}_{t-k}\|_q\le L_p\eta\sum_{j=t-k}^{t-1}\|\mathbf{u}_j\|_p$.
Minkowski's inequality in $L^2$ and~\eqref{eq:update-moment} therefore give
\begin{equation}
(\mathbb E\|\mathbf{r}_t\|_q^2)^{1/2}
\le G\gamma^t+hL_p\eta\sqrt M\sum_{k=1}^{t-1}k\gamma^k
\le G\gamma^t+\frac{L_p\eta\sqrt M\gamma}{h}.
\end{equation}
Also $\|\bar{\mathbf{g}}_t\|_q\le G$ and $h\|\bar{\mathbf{G}}_t\|_q\le G(1-\gamma^t)$, proving the pointwise $2G$ bound. Cauchy--Schwarz gives the first product bound without assuming independence. The second uses $\|\mathbf{r}_t\|_q^q\le(2G)^{q-1}\|\mathbf{r}_t\|_q$ followed by $\mathbb E\|\mathbf r_t\|_q\le(\mathbb E\|\mathbf r_t\|_q^2)^{1/2}\le R_t^{\rm drift}$. In particular, we never substitute the moment bound $M$ as a pointwise bound on $\|\mathbf{u}_t\|_p^2$.
\end{proof}

\begin{lemma}[Alignment with the current gradient]\label{lemma:alignment}
Let $A,B$ be defined in~\eqref{eq:theory-constants}. Then
\begin{equation}\label{eq:alignment-correct}
\mathbb E\langle\bar{\mathbf{g}}_t,\mathbf{u}_t\rangle
\ge A\mathbb E\|\bar{\mathbf{g}}_t\|_q^q
-hK_\beta\mathbb E\|\mathbf{N}_t\|_q^q-BR_t^{\rm drift}.
\end{equation}
\end{lemma}
\begin{proof}
The scalar perturbation bound followed by~\eqref{eq:young-explicit} gives
\begin{align}
\langle\bar{\mathbf{G}}_t,\Phi_\beta(\bar{\mathbf{G}}_t+\mathbf{N}_t)\rangle
&\ge\|\bar{\mathbf{G}}_t\|_q^q-2^{1-\beta}\sum_i|\bar G_{t,i}||N_{t,i}|^\beta\nonumber\\
&\ge\tfrac12\|\bar{\mathbf{G}}_t\|_q^q-K_\beta\|\mathbf{N}_t\|_q^q.
\end{align}
Since $\bar{\mathbf{g}}_t=h\bar{\mathbf{G}}_t+\mathbf{r}_t$, dual-pair H\"older and Lemma~\ref{lemma:drift} imply
\begin{equation}
\mathbb E\langle\bar{\mathbf{g}}_t,\mathbf{u}_t\rangle
\ge\tfrac h2\mathbb E\|\bar{\mathbf{G}}_t\|_q^q-hK_\beta\mathbb E\|\mathbf{N}_t\|_q^q-\sqrt M R_t^{\rm drift}.
\end{equation}
Finally $\|\bar{\mathbf{g}}_t\|_q^q\le2^{q-1}(h^q\|\bar{\mathbf{G}}_t\|_q^q+\|\mathbf{r}_t\|_q^q)$, hence
\begin{equation}
\tfrac h2\mathbb E\|\bar{\mathbf{G}}_t\|_q^q
\ge2^{-q}h^{1-q}\mathbb E\|\bar{\mathbf{g}}_t\|_q^q
-\tfrac12h^{1-q}(2G)^{q-1}R_t^{\rm drift}.
\end{equation}
Substituting $q-1=\beta$ proves the result.
\end{proof}

Integrating the gradient Lipschitz assumption along a line segment gives
$f(\mathbf{y})\le f(\mathbf{x})+\langle\nabla f(\mathbf{x}),\mathbf{y}-\mathbf{x}\rangle+(L_p/2)\|\mathbf{y}-\mathbf{x}\|_p^2$.
The update and the moment bound imply
\begin{equation}
\mathbb E f(\bm{\theta}_t)\le\mathbb E f(\bm{\theta}_{t-1})
-\eta\mathbb E\langle\bar{\mathbf{g}}_t,\mathbf{u}_t\rangle+\tfrac{L_pM}2\eta^2.
\end{equation}
The bounded-gradient assumption and finite update moments ensure integrability at each finite time. Summing and using $f\ge f^*$ yields
\begin{equation}
\eta\sum_{t=1}^T\mathbb E\langle\bar{\mathbf{g}}_t,\mathbf{u}_t\rangle\le\Delta_0+\tfrac{L_pM}2T\eta^2.
\end{equation}
Insert~\eqref{eq:alignment-correct}, use~\eqref{eq:noise-moment}, and note $\sum_{t=1}^TR_t^{\rm drift}\le TD\eta+G\gamma/h$:
\begin{equation}\label{eq:master-bound}
A\frac1T\sum_{t=1}^T\mathbb E\|\bar{\mathbf{g}}_t\|_q^q
\le\frac{\Delta_0}{\eta T}+\left(\tfrac{L_pM}2+BD\right)\eta
+\frac{BG\gamma}{hT}+K_\beta h^{1-q}\frac1T\sum_{t=1}^T\nu_t.
\end{equation}
Since $h^{1-q}/A=2^q$, substituting $\eta=c/\sqrt T$ gives~\eqref{eq:corrected-rate}. For the squared Euclidean gradient, $|\bar g_{t,i}|\le G$ implies pointwise
\begin{equation}\label{eq:correct-norm-conversion}
\|\bar{\mathbf{g}}_t\|_2^2=\sum_i|\bar g_{t,i}|^2
\le G^{2-q}\sum_i|\bar g_{t,i}|^q=G^{1-\beta}\|\bar{\mathbf{g}}_t\|_q^q.
\end{equation}
Take expectations and use that the minimum is at most the average. This proves the theorem.\hfill$\square$

\subsection{Conditions That Remove the Residual}
\begin{corollary}[Deterministic and vanishing noise]\label{cor:vanishing}
In Theorem~\ref{thm:convergence}, $\sigma=0$ gives $O(T^{-1/2})$ for the average squared gradient. More generally, if $\nu_t\le\bar\nu_t\le\sigma^q$ for a fixed deterministic sequence with $T^{-1}\sum_{t=1}^T\bar\nu_t\to0$, the bound tends to zero as the horizon grows. This is a guarantee for a family of horizon-dependent runs, not an almost-sure convergence claim for a single infinite run.
\end{corollary}
\begin{proof}
Use the actual moments $\nu_t$ in~\eqref{eq:corrected-rate}; the uniform envelope keeps $M,D,B$ fixed. Both terms on the right tend to zero.
\end{proof}

\begin{corollary}[Relative-noise condition]\label{cor:growth}
Replace the noise assumption by
$\mathbb E[\|\bm{\xi}_t\|_q^q\mid\mathcal F_{t-1}]\le\rho^q\|\bar{\mathbf{g}}_t\|_q^q$.
If $2^qK_\beta\rho^q\le1/2$, then
\begin{equation}
\frac1T\sum_{t=1}^T\mathbb E\|\bar{\mathbf{g}}_t\|_q^q\le2R_T=O(T^{-1/2}),
\end{equation}
where $R_T$ uses $\sigma=\rho G$. The same rate holds for the average squared Euclidean gradient after multiplication by $G^{1-\beta}$.
\end{corollary}
\begin{proof}
Bounded gradients imply the uniform envelope $\sigma=\rho G$, so every lemma still applies. In~\eqref{eq:corrected-rate}, substitute $\nu_t\le\rho^q\mathbb E\|\bar{\mathbf{g}}_t\|_q^q$ and move this term to the left. The assumed coefficient is at most $1/2$.
\end{proof}
This is an explicit, restrictive sufficient condition, especially for small $\beta$. It is not implied by unbiasedness or ordinary bounded variance, and it has not been checked in our experiments.

\begin{corollary}[Growing independent mini-batches]\label{cor:batch}
Suppose the oracle averages $b_t$ conditionally independent, unbiased per-sample gradients with conditional Euclidean noise variance at most $v^2$. In addition to Assumption~\ref{assum:regularity}, this gives
\begin{equation}
\nu_t\le d^{1-q/2}v^q b_t^{-q/2}.
\end{equation}
Any $b_t\to\infty$ removes the residual in the horizon limit. If $b_t=\lceil a t\rceil$ for $a>0$, the squared-gradient bound is $O(T^{-1/2})$ for every fixed $\beta>0$.
\end{corollary}
\begin{proof}
Conditional independence gives $\mathbb E[\|\bm{\xi}_t\|_2^2\mid\mathcal F_{t-1}]\le v^2/b_t$. The norm inequality $\|\mathbf{x}\|_q\le d^{1/q-1/2}\|\mathbf{x}\|_2$ and concavity of $z^{q/2}$ give the stated moment bound. Ces\`aro averaging proves the first assertion. For $1<q<2$, $T^{-1}\sum_{t=1}^Tt^{-q/2}=O(T^{-q/2})=O(T^{-1/2})$; for $q=2$, it is $O(\log(T)/T)=O(T^{-1/2})$. Apply Theorem~\ref{thm:convergence} with the uniform envelope $\sigma^q=d^{1-q/2}v^q$.
\end{proof}
A mere finite $q$-moment does not imply the $b_t^{-q/2}$ scaling; the per-sample second-moment assumption above supplies it. Linear batch growth uses $N=\Theta(T^2)$ stochastic gradient evaluations, yielding only $O(N^{-1/4})$ for the squared-gradient bound derived here. We do not call this sample-optimal.

\subsection{Interpretation and Limits}
For $q\le2\le p$, norm equivalence gives the valid conversions in~\eqref{eq:dimension-conversion}: applying $\|\mathbf{z}\|_q\le d^{1/q-1/2}\|\mathbf{z}\|_2$ to gradient differences and $\|\mathbf{x}-\mathbf{y}\|_2\le d^{1/2-1/p}\|\mathbf{x}-\mathbf{y}\|_p$ to displacements yields $L_p=d^{1/q-1/p}L_2$. The bounded-gradient and noise conversions follow in the same way, with Jensen for the noise. Thus native notation removes explicit conversion factors from the proof, but does not remove possible dimensional scaling of the assumptions.

The residual is an upper-bound term, not a matching lower bound or a claim that every trajectory stalls. Unbiased noise alone does not commute with the nonlinear transform: for example, a scalar noise taking $2$ with probability $1/3$ and $-1$ with probability $2/3$ has zero mean but $\mathbb E\Phi_\beta(\xi)=(2^\beta-2)/3\ne0$ for $\beta<1$. Additional distributional or variance-reduction assumptions therefore matter. This example explains why exact alignment cannot simply be inferred from unbiasedness; it is not a counterexample to every fixed-momentum run.

Our proof uses geometric-sum convexity, not a second-moment martingale identity for $q<2$. It also retains the initialization term $G\gamma^t$. It supplies neither a guarantee for quantization, clipping, and decoupled weight decay nor a proof for scheduled momentum or the practical cosine schedule. Those extensions require separate analyses.


\end{document}

\end{document}
